\subsection{The glued group}

We first realize the incidence data by finite subgroups of one group.
For each $v\in V$, let $K_v^0$ be a coordinatized copy of the additive
group of $\F_\ell^2$, and for each $f\in\mathcal F$, let $L_f^0$ be a
copy of the additive group of $\F_\ell$.  Since exactly $\ell+1$
hyperedges contain $v$, choose injections
\[
  \iota_{vf}:L_f^0\longrightarrow K_v^0\qquad(v\in f)
\]
whose images are the $\ell+1$ distinct one-dimensional subspaces of
$K_v^0$.  These are homomorphisms of additive groups.  In particular,
their images partition $K_v^0\setminus\{0\}$ outside their common
zero element.

Introduce a generator $s_{f,c}$ for each $f\in\mathcal F$ and each
$0\ne c\in L_f^0$.  A signed word
\[
  s_{f_1,c_1}^{\varepsilon_1}\cdots
  s_{f_k,c_k}^{\varepsilon_k},\qquad \varepsilon_i\in\{1,-1\},
\]
is a defining relator whenever some $v\in V$ belongs to every $f_i$
and
\[
  \sum_{i=1}^k\varepsilon_i\iota_{vf_i}(c_i)=0
  \quad\text{in }K_v^0.
\]
Let $N$ be the group with these generators and relators.  It is
finitely generated, since the incidence data and all the local groups
are finite.

There is a homomorphism $\theta_v:K_v^0\to N$ determined by
\[
  \theta_v(0)=1,\qquad
  \theta_v\bigl(\iota_{vf}(c)\bigr)=s_{f,c}\quad(c\ne0).
\]
Indeed, each nonzero element has a unique expression on the left.
For $a,b\in K_v^0$, the word expressing
$\theta_v(a)\theta_v(b)\theta_v(a+b)^{-1}$ is a defining relator,
after omitting factors corresponding to zero elements.

The same group is presented by the finite sublist of defining relators
of length at most three. These include every local multiplication-table
relation just displayed, as well as the local inverse relations. Within
one $K_v^0$, successive applications of those relations reduce any
signed word to the generator for its sum, or to the empty word when
the sum is zero. Thus they imply every defining relator of $N$, and
$N$ is finitely presented.
We retain the all-word formulation because it makes the planar proof
below more transparent.

\begin{lemma}\label{lem:embedding}
If the incidence graph of $(V,\mathcal F)$ has girth at least $12$,
then every homomorphism $\theta_v:K_v^0\to N$ is injective.
\end{lemma}

\begin{proof}
Every nonzero element of $K_v^0$ maps to a single generator $s_{f,c}$.
It therefore suffices to prove that no such generator is the identity
in $N$.  Suppose, to the contrary, that $s_{f,c}=1$, with $c\ne0$.
We will encode this equality by a planar graph with zero local sums
at every vertex except one.  Reducing the graph will make the local
group relations force lower bounds on vertex degrees, while incidence
girth forces lower bounds on facial lengths.  These bounds will
contradict Euler's formula even with the one exceptional vertex.

\emph{From a word equality to a planar graph.}
The one-letter word $s_{f,c}$ can be changed to the empty word by
finitely many insertions and deletions of contiguous defining
relators.  To justify this form of reduction, take the congruence on
signed words generated by those operations.  Every formal inverse
pair is a defining relator, since its hyperedge has an incident
vertex.  The quotient is consequently a group, and its universal
property is exactly that of the presentation of $N$.  Thus equality
in $N$ is equality under this congruence.  Equivalently, each factor
$u r^{\pm1}u^{-1}$ in a finite normal-closure expression can be
inserted by first forming $u u^{-1}$ through nested inverse-pair
insertions and then inserting $r^{\pm1}$ between its two halves.
Inverses of defining relators are again defining relators.

Draw the successive words on horizontal levels, treating each signed
letter as an individual token.  Between two levels, connect the
tokens that persist in their unchanged order.  An inserted or deleted
block begins or ends its tokens at a new black vertex, labelled by a
point $v$ at which the corresponding relator has zero sum.  This can
be drawn without crossings: the affected block is contiguous, so its
star occupies a single gap between the other order-preserving wires.
Joining the finitely many strips gives a planar drawing.

Every inserted token $s_{i,y}^{\varepsilon}$ has one creation vertex
and one deletion vertex.  Subdivide its full lifetime wire by a white
vertex labelled $i$, and give its two edges weights $\varepsilon y$
at the creation end and $-\varepsilon y$ at the deletion end.  The
initial token instead has an exceptional white endpoint labelled $f$,
joined to its deletion vertex by an edge of weight $-c$.  We obtain
a finite planar bipartite multigraph with black labels in $V$ and
white labels in $\mathcal F$.  An edge joining labels $v$ and $i$
has a nonzero weight in $L_i^0$, and $v\in i$.  Its balances are
\begin{equation}\label{eq:embedding-balances}
\begin{aligned}
  &\sum_{a\ni x}\iota_{v,i(a)}\bigl(\omega(a)\bigr)=0
    &&\text{at each black vertex $x$ labelled $v$},\\
  &\sum_{a\ni y}\omega(a)=0
    &&\text{at each ordinary white vertex $y$},
\end{aligned}
\end{equation}
where $\omega(a)$ is the weight of edge $a$ and $i(a)$ is its white
label.  The exceptional white vertex has nonzero sum.  At creation
vertices the first balance is the relator sum, and at deletion
vertices it is its negative.  Ordinary white sums vanish because
their two weights are opposite.  Retain only the connected component
of the exception.

\emph{Reducing the planar graph.}
Among all connected planar bipartite multigraphs with the indicated
labels, nonzero edge weights, and balances
\eqref{eq:embedding-balances}, and with exactly one exceptional white
vertex of nonzero sum, choose one minimizing lexicographically the
number of vertices and then the number of edges.  We do not require
this comparison graph to come from a word reduction.  Its exceptional
vertex has positive degree, so the graph has an edge.

There are no parallel edges.  Two parallel edges have weights in the
same $L_i^0$ and the same map $\iota_{vi}$ at their black endpoint.
Replace them by one edge carrying their sum, or delete both if that
sum is zero, and again retain the component of the exception.
Additivity preserves every balance and preserves the exceptional
nonzero sum.  This gives a smaller comparison graph, a contradiction.

Fix a spherical embedding of the minimal graph.  We claim that, at
each vertex of degree at least two, cyclically consecutive neighbors
have distinct labels.  If two such neighbors have the same label,
they are distinct vertices because parallel edges have been removed.
Join them by an auxiliary arc following one side of the two incident
edges and the empty sector between those edges, and contract the
arc.  A sufficiently narrow neighborhood of this two-edge path
makes the arc disjoint from the rest of the graph.  This construction
is local and remains valid when either edge is a bridge.

The contraction merges vertices of the same color and label.  Each
original edge keeps its weight and its opposite-color endpoint; hence
the resulting graph is still bipartite, and every edge weight is
nonzero.  Parallel edges created by the contraction are allowed in
the comparison class.  When black vertices are merged, their
balances add in the same $K_v^0$, while white balances are unchanged.
When white vertices are merged, their balances add in the same
$L_i^0$, while every black incidence map and weighted sum is
unchanged.  If the exception is merged, its sum is increased by zero.
Thus the contraction gives a comparison graph with fewer vertices,
proving the claim.

\begin{figure}[htbp]
  \centering
  \begin{tikzpicture}[
    x=1cm, y=1cm, font=\small,
    merge white/.style={circle, draw, fill=white, line width=.6pt,
      minimum size=5.5pt, inner sep=0pt},
    merge black/.style={circle, draw, fill=black,
      minimum size=4.5pt, inner sep=0pt},
    merge edge/.style={line width=.6pt},
    merge auxiliary/.style={densely dashed, line width=.55pt}
  ]
    \begin{scope}
      \coordinate (merge-before-p-position) at (-1.1,1.65);
      \coordinate (merge-before-q-position) at (1.1,1.65);
      \coordinate (merge-before-x-position) at (0,0);
      \draw[merge edge]
        (-1.55,2.2) -- (merge-before-p-position)
        -- (merge-before-x-position) -- (merge-before-q-position)
        -- (1.55,2.2);
      \draw[merge auxiliary]
        (merge-before-p-position)
        .. controls (-.92,1.47) and (-.15,.43) .. (0,.43)
        .. controls (.15,.43) and (.92,1.47) ..
        (merge-before-q-position);
      \node[merge white, label=above:{$f$}, label=left:{$p$}]
        (merge-before-p) at (merge-before-p-position) {};
      \node[merge white, label=above:{$f$}, label=right:{$q$}]
        (merge-before-q) at (merge-before-q-position) {};
      \node[merge black, label=below:{$v$}, label=left:{$x$}]
        (merge-before-x) at (merge-before-x-position) {};
      \node at (-.76,.78) {$a$};
      \node at (.76,.78) {$b$};
    \end{scope}

    \draw[-{Stealth[length=5pt]}, line width=.6pt]
      (1.95,1) -- node[above=3pt] {contract} (3.35,1);

    \begin{scope}[xshift=5.1cm]
      \coordinate (merge-after-r-position) at (0,1.65);
      \coordinate (merge-after-x-position) at (0,0);
      \draw[merge edge]
        (-.75,2.2) -- (merge-after-r-position) -- (.75,2.2);
      \draw[merge edge]
        (merge-after-x-position)
        .. controls (-.63,.45) and (-.63,1.2) ..
        (merge-after-r-position);
      \draw[merge edge]
        (merge-after-x-position)
        .. controls (.63,.45) and (.63,1.2) ..
        (merge-after-r-position);
      \node[merge white, label=above:{$f$}, label=right:{$r$}]
        (merge-after-r) at (merge-after-r-position) {};
      \node[merge black, label=below:{$v$}, label=left:{$x$}]
        (merge-after-x) at (merge-after-x-position) {};
      \node at (-.72,.82) {$a$};
      \node at (.72,.82) {$b$};
    \end{scope}
  \end{tikzpicture}
  \caption{The white vertices $p,q$, both labelled $f$, merge to $r$
    along the dashed auxiliary arc in the empty sector between their
    consecutive edges at the black vertex $x$, labelled $v$.
    The edges of weights $a,b$ remain distinct parallel edges.}
  \label{fig:merge}
\end{figure}


\emph{Degrees, faces, and Euler's formula.}
An ordinary white vertex cannot have degree one, since its incident
weight would be nonzero.  A black vertex cannot have degree one
either, since the maps $\iota_{vi}$ are injective.  If a black vertex
had degree two, its two white neighbors would have distinct labels
by the preceding claim.  Its zero balance would then equate two
nonzero vectors in distinct one-dimensional subspaces of $K_v^0$,
which is impossible.  Consequently every black degree is at least
three, and every white degree is at least two except possibly the
exceptional degree, which is at least one.

Consider the facial boundary walks of this connected plane graph,
counting edge sides with multiplicity.  At a turn through a vertex
of degree at least two, the entering and leaving edges are distinct
and consecutive in its rotation.  Their opposite endpoints have
different labels.  Thus the image of the boundary walk in the
incidence graph has no immediate backtracking at that turn.  This
also applies to the cyclic closing turn of the walk.  Bridges cause
no difficulty: their two sides are counted separately, and a facial
walk reverses immediately along an edge only at a degree-one vertex.
There is at most one such vertex, the exceptional white vertex, and
it belongs to just one face.

It follows that all but at most one face have boundary walks whose
images are nonempty closed walks without immediate backtracking in
the incidence graph.  Such a walk contains a cycle, so each of these
faces has length at least $12$.  Write $n_{\rm b},n_{\rm w},n_{\rm f}$
for the numbers of black vertices, white vertices, and faces, and
$E$ for the number of edges.  The degree bounds and the total of
$2E$ edge sides give
\[
  n_{\rm b}\le\frac E3,\qquad
  n_{\rm w}\le\frac{E+1}{2},\qquad
  n_{\rm f}\le\frac E6+1.
\]
Therefore
\[
  n_{\rm b}+n_{\rm w}+n_{\rm f}
  \le E+\frac32<E+2,
\]
contrary to Euler's formula for a connected plane graph.  The
assumed word equality is impossible, and every $\theta_v$ is
injective.
\end{proof}

We henceforth identify $K_v^0$ with its image $K_v\le N$, and identify
$L_f^0$ with the subgroup $L_f\le N$ given by $c\mapsto s_{f,c}$,
with $0$ mapped to the identity.  The latter is injective by choosing
any point incident to $f$ and applying Lemma~\ref{lem:embedding}.
The shared generators identify this same subgroup inside every
$K_v$ with $v\in f$.  In particular, $|K_v|=\ell^2$, $|L_f|=\ell$,
and the incident subgroups $L_f$ partition $K_v$ outside the
identity.  No assertion about intersections between other local
subgroups is needed.
